\chapter{Conclusion and Future Work}
\label{ch:conclusion}

This report has documented \tool, a static timing analysis and
logical-effort-guided gate-sizing tool built to the specification of the
Advanced VLSI Design course project, together with two substantial
additions beyond that specification: a synthetic benchmark generation and
oracle-based evaluation framework, and an interactive topology viewer. The
core analysis pipeline --- structural validation, DAG construction,
RC-based rise/fall-separated delay modeling with unate-aware timing arcs,
forward and backward static timing analysis, logical-effort path analysis,
power and area estimation, and Monte Carlo statistical timing analysis ---
was traced in this report down to the specific class and function names
that implement each stage.

On the fixed benchmark suite, the logical-effort-guided canonical
heuristic resolved $13$ of $14$ originally-violating circuits, compared to
$7$ of $14$ for a logical-effort-free random-greedy baseline, while
simultaneously issuing fewer total STA calls and reaching a lower final
cost on thirteen of fifteen circuits (\cref{ch:results}). Two case studies,
traced directly against the recorded optimizer iteration logs rather than
inferred from final outcomes alone, showed that this advantage comes from
a better-informed \emph{search order} interacting favorably with the
area/power budget, not from any difference in what a given heuristic is
permitted to consider. The one Monte Carlo--instrumented circuit
demonstrated that a nominal-corner timing improvement from optimization
corresponded to a real, substantial improvement in statistical timing
yield (from 36.9\% to 82.6\%), and that the sizing changes shifted which
path is critical under process variation, not merely how much margin the
original critical path retained.
